\documentclass[
    ngerman,
    accentcolor=3b,
    dark_mode,
    fontsize= 12pt,
    a4paper,
    aspectratio=169,
    colorback=true,
    fancy_row_colors,
    leqno,
    fleqn,
    boxarc=3pt,
    fleqn,
    main,
    % shell_escape = false, % Kompatibilität mit sharelatex
    ]{algoslides}

%%------------%%
%%--Packages--%%
%%------------%%

% \usepackage{audutils}
% \usepackage{fopbot}
\usepackage{booktabs}
\usepackage{tipa}


%%----------------------------%%
%%--Stilistische Anpassungen--%%
%%----------------------------%%


\renewcommand\tabularxcolumn[1]{m{#1}}% for vertical centering text in X column
% Remove unwanted space from tables
\aboverulesep = 0mm \belowrulesep = 0mm
\renewcommand{\arraystretch}{1.4}

%%---------------------------%%
%%--Dokumenteneinstellungen--%%
%%---------------------------%%

% Workshop-Individuelle Einstellungen
\def\workshoptitle{\LaTeX}
\def\shortworkshoptitle{\LaTeX}
\subtitle{Einsteiger Workshop}
\def\gruesswoerte{Gude}
\author{Ruben Deisenroth}
\titlegraphic*{\includegraphics{LaTeX-Wallpaper.png}}


% Ofahrt-Einstellungen
\graphicspath{{pictures}}
\title[\shortworkshoptitle{}]{\workshoptitle{} -- Ofahrt 2022}
\department{TU Darmstadt | Fachbereich Informatik | Ofahrt 2022 | \insertshorttitle{}}
\date{\today}
\logo*{\includegraphics{example-image-16x9}} % TODO: Ofahrt-Logo erstellen
\def\codeDir{code}
%%-------------------------%%
%%--Beginn des Dokumentes--%%
%%-------------------------%%

\begin{document}

    %%-----------%%
    %%--Titelei--%%
    %%-----------%%

    \maketitle{}

    \begin{frame}[c]
        % Die Begrüßenden Worte können individuell pro Workshop festgelegt werden. Am Besten nur was kurzes, sowie "Gude", "Hi", "Herzlich Willkommen!", ...
        \centering\huge\textbf{\gruesswoerte{}}
    \end{frame}

    %%---------------------------%%
    %%--Beginn der Präsentation--%%
    %%---------------------------%%
    
    \begin{frame}
        \frametitle{Das steht heute auf dem Plan}
        \tableofcontents[subsubsectionstyle=hide]
    \end{frame}


    %%-----------%%
    %%--Kapitel--%%
    %%-----------%%

    \input{chapters/chapter01}
    \input{chapters/chapter02}
    \documentclass[
    ngerman,
    accentcolor=3b,
    dark_mode,
    fontsize= 12pt,
    a4paper,
    aspectratio=169,
    colorback=true,
    fancy_row_colors,
    leqno,
    fleqn,
    boxarc=3pt,
    fleqn,
    % shell_escape = false, % Kompatibilität mit sharelatex
]{algoslides}

%%------------%%
%%--Packages--%%
%%------------%%

% \usepackage{audutils}
% \usepackage{fopbot}
\usepackage{booktabs}
\usepackage{tipa}

% Import all Packages from Main Preamble with relative Path (buggy, list packages instead)
% \subimport*{../../}{preamble}

%%--------------------------%%
%%--Imports from Main File--%%
%%--------------------------%%

% Get Labels from Main Document using the xr-hyper Package
\externaldocument[ext:]{../main}
% Set Graphics Path, so pictures load correctly
\graphicspath{{../pictures}}
\def\codeDir{../code}

\begin{document}
    \section{\LaTeX-Basics}\label{2}\label{LaTeX-Basics}
    \subsection{Syntax}
    \begin{frame}[fragile]
        \slidehead{}
        \begin{itemize}
            \item Kommandos beginnen mit einem Backslask: \verb+\+
            \item Kommandonamen enthalten keine Zahlen, Leerzeichen und Sonderzeichen
            \item Argumente werden in geschweiften Klammern angegeben \verb+{}+
            \item Optionale Argumente werden in eckigen Klammern angegeben \verb+[]+
            \item Environments werden mit \verb+\begin{<name>}+ und \verb+\end{<name>}+ angegeben
            \item Kommentare werden mit einem \verb+%+ beginnen
        \end{itemize}
    \end{frame}
    \subsection{Aufbau eines Dokumentes}\label{2.1}\label{2.1}
    \begin{frame}[fragile]
        \slidehead{}
        \begin{columns}[c]
            \begin{column}{.5\textwidth}
                \inputCode[]{minted language=latex,title=\codeBlockTitle{Minimales \LaTeX-Dokument}}{\codeDir/minimal.tex}
            \end{column}%
            \begin{column}{.5\textwidth}
                \centering
                \colorbox{white}{%
                    \begingroup
                    \color{black}
                    \includegraphics[width=\textwidth,height=5cm,keepaspectratio]{\codeDir/minimal}
                    \endgroup
                }
            \end{column}
        \end{columns}
    \end{frame}
    \begin{frame}[fragile]
        \slidehead{}
        \begin{columns}[c]
            \begin{column}{.5\textwidth}
                \inputCode[]{minted language=latex,title=\codeBlockTitle{Minimales \LaTeX-Dokument}}{\codeDir/minimal-documented.tex}
            \end{column}%
            \begin{column}{.5\textwidth}
                \centering
                \colorbox{white}{%
                    \begingroup
                    \color{black}
                    \includegraphics[width=\textwidth,height=5cm,keepaspectratio]{\codeDir/minimal-documented}
                    \endgroup
                }
            \end{column}
        \end{columns}
    \end{frame}
\end{document}
\end{document}
